% !TEX root = ../my-thesis.tex
%
\chapter{Appendix}\label{sec:appendix}

% Style für prompt Darstellung
\lstdefinestyle{promptstyle}{
  basicstyle=\ttfamily\footnotesize,
  columns=fullflexible,
  keepspaces=true,
  upquote=true,
  showstringspaces=false,
  breaklines=true,
  breakindent=0pt,
  breakautoindent=false,
  postbreak=\mbox{\textcolor{gray}{$\hookrightarrow$}\space},
  frame=single,
  rulecolor=\color{gray!50},
  xleftmargin=0pt,
  numbers=none,
  captionpos=b,
}

\section{Single-Shot Prompt}\label{sec:appendix:single-shot-prompt}

\begin{lstlisting}[style=promptstyle,
                   caption={Ausgabeprotokoll der Single-Shot-Solution
                            (Python-Konstante \texttt{FORMAT\_INSTRUCTION});
                            wird der Task-Beschreibung unverändert angehängt.},
                   label={lst:prompt-singleshot}]
Write out the complete project, one file at a time. Introduce each file with a
line naming its path, then give that file's entire contents in a fenced code
block:

### FILE: src/example/core.py
```python
def add(left: int, right: int) -> int:
    return left + right
```

- Paths are relative to the project root: `pyproject.toml` belongs at
  `pyproject.toml`, not inside a directory named after the project.
- Where the specification shows a directory tree, its top-level directory is
  the project root itself and not a directory to create.
- No absolute paths, and no `..`.
- Include every file the project needs in order to be installed and used,
  packaging and configuration among them.
- Write each file out in full. Do not abbreviate a body to a comment, and do
  not leave a placeholder to be filled in later.
- If a file's own contents contain a fenced code block, fence that file with
  more backticks than appear anywhere inside it.
- Only the file blocks are read. Any other text is ignored.
\end{lstlisting}